\setcounter{page}{1}
\chapter*{Resumen}

Este trabajo describe el desarrollo e implementación de un protocolo integral de
calibración y optimización para un sistema de microtomografía computarizada de rayos X ($\mu$CT)
de arquitectura abierta. El objetivo central fue transformar un conjunto de componentes
independientes en un instrumento de alta precisión para la investigación científica,
comparable a un dispositivo comercial.

La metodología comenzó con la estabilización del hardware, fijando la fuente y el
detector sobre una base rígida para eliminar desalineaciones accidentales. Se
implementaron técnicas de alineación física mediante el análisis de proyecciones
cónicas y la superposición de proyecciones de un cilindro, logrando reducir desviaciones angulares
a menos de 5$^{\circ}$ y aproximando el punto principal del haz al centro del detector. En
la etapa de calibración analítica, se evaluaron dos métodos: uno
de múltiples proyecciones, estudiando las trayectorias elípticas de esferas metálicas,
que resultó ser el más robusto y consistente para
determinar parámetros geométricos finos; y el de una única proyección, a partir de
la relación entre los lados de un cuadrilátero proyectado, el cual se
descartó por su extrema sensibilidad a errores de posicionamiento.

Asimismo, se caracterizó la resolución espacial efectiva del sistema a partir de un
estudio de la función de transferencia de modulación
(MTF). Los resultados situaron la resolución entre 15 y 60 $\mu$m, y se obtuvo la
curva empírica en función de la magnificación. Finalmente, se aplicaron técnicas de pre-procesamiento
digital como la corrección de campo plano (FFC) y la corrección de uniformidad
temporal, las cuales eliminaron patrones electrónicos del detector, aumentaron
la relación señal-ruido (SNR) y uniformizaron la intensidad del haz en área y tiempo
de un conjunto de proyecciones. El protocolo fue validado
exitosamente mediante la reconstrucción tridimensional de muestras biológicas y
mecánicas con alta nitidez y ausencia de artefactos, obteniéndose un sistema de microtomografía
con resultados equivalentes a los de uno comercial.
